\section*{Chapitre \ref{ch:rs:portfolio-construction-i} --- Construction de portefeuille I}

\begin{solution}{exo:rs:portfolio-construction-i:1}
$w_i = \alpha_i/(\gamma\sigma_i^2)$~: $0.02/(5 \times 0.04) = 0.1$ et $0.01/(5 \times 0.01) = 0.2$. $\mathrm{IR} = \sqrt{0.02^2/0.04 + 0.01^2/0.01} = \sqrt{0.02} = 0.14$.
\end{solution}

\begin{solution}{exo:rs:portfolio-construction-i:2}
$84 \times 12 \times 0.0010 = 1.01$~: environ la totalité du capital du portefeuille chaque année, ce qui explique que son ratio d'information passe de 0.92 à 0.43.
\end{solution}

\begin{solution}{exo:rs:portfolio-construction-i:3}
$\sqrt{(1 + \mathrm{IR}^2/2)/T}$~: $\sqrt{(1 + 1.77^2/2)/7.9} = 0.57$ et $\sqrt{(1 + 1.39^2/2)/7.9} = 0.50$. Les deux valeurs diffèrent de moins d'une erreur type.
\end{solution}

\begin{solution}{exo:rs:portfolio-construction-i:4}
Avec 126 observations de 100 rendements, les plus petites valeurs propres de la covariance empirique sont très inférieures aux vraies (au rapport 100/126, beaucoup sont proches de zéro). Le
portefeuille optimal $\Sigma^{-1}\alpha/\gamma$ divise par elles~: les combinaisons presque sans risque \omterm{def:rs:overfitting:overfitting}{dans l'échantillon} reçoivent des poids énormes, longs comme courts, et l'exposition brute explose.
Leur risque réel n'est pas faible, si bien que la volatilité réalisée (207~\%) écrase la prévision (36~\%).
\end{solution}

\begin{solution}{exo:rs:portfolio-construction-i:5}
La prévision est un signal vrai et persistant, auquel s'ajoute un bruit nouveau chaque mois. Un portefeuille qui ne rejoint la cible de chaque mois que partiellement détient une moyenne de
des cibles récentes, dans laquelle le bruit se compense en partie et le signal non. Avant frais, il échange un peu moins de signal contre beaucoup moins de bruit.
\end{solution}

\begin{solution}{exo:rs:portfolio-construction-i:6}
$0.05 \times 10\,\mathrm{m}/1\,000\,\mathrm{m} = 0.05~\%$ du capital.
\end{solution}

\begin{solution}{exo:rs:portfolio-construction-i:7}
La limite brute de 2 n'est saturée que si le portefeuille la dépasserait autrement. Avec les limites de liquidité et de rotation, l'exposition brute vaut en moyenne 1.01, si bien que la limite est lâche chaque
mois~: son multiplicateur est nul, de même que son effet et sa part (conditions de complémentarité).
\end{solution}

\begin{solution}{exo:rs:portfolio-construction-i:8}
Le ratio ex ante mesure les erreurs propres de l'optimiseur~: avec une covariance empirique sur six mois portant sur une centaine de titres, la volatilité réalisée du portefeuille atteignait 207~\% contre 36~\%
prévue, et son ratio réalisé vaut 0.92 avant frais et 0.43 après. Il faut dimensionner sur la performance réalisée, nette des coûts, d'un portefeuille construit avec un modèle factoriel de risque et des
alphas mis à l'échelle, avec son erreur type.
\end{solution}

\begin{solution}{pb:rs:portfolio-construction-i:1}
\begin{enumerate}
\item $w^* = \Sigma^{-1}\alpha/\gamma$~; $\mathrm{IR} = \sqrt{\alpha^\top\Sigma^{-1}\alpha}$.
\item L'alpha est le rendement mensuel attendu propre à chaque action de la simulation, auquel s'ajoute un bruit deux fois plus grand, standardisé et mis à l'échelle par $\mathrm{IC}\,\sigma_i$ avec
$\mathrm{IC} = 0.05$~; le modèle de risque est celui du chapitre~24, à horizon mensuel. La vérité augmentée d'un bruit isole la construction de la prévision.
\item Exposition brute 61~; volatilité prévue 36~\%, volatilité réalisée 207~\%~; ratio ex ante 4.25, ratio réalisé 0.92 et 0.43 après frais~; rotation de 84 par mois.
\item Les plus petites valeurs propres de la covariance empirique correspondent à des combinaisons de nombreux titres~; les deux premières positions représentent en moyenne 9~\% de l'exposition brute.
\item La volatilité prévue et la volatilité réalisée concordent (13.7~\% et 14.0~\%), et le ratio d'information net vaut 1.39.
\item De 1.64 à 1.59~: l'alpha est propre à chaque action, si bien que le portefeuille neutre perd peu.
\item L'exposition brute baisse de 4.35 à 2.21 et 2.00, et les frais de 6.8~\% à 3.0~\% et 2.8~\% par an.
\item 28 en moyenne~; le ratio ex ante baisse de 1.44 à 1.23.
\item Les ratios réalisés sur 7.9 ans ont des erreurs types proches de 0.5.
\item Rotation de 0.53 par mois, frais de 0.6~\% par an.
\item Il calcule la moyenne des prévisions mensuelles, dont le bruit est indépendant d'un mois sur l'autre.
\item Le portefeuille ne suit plus la prévision bruitée~; sa corrélation avec elle baisse (de 0.97 à 0.58) tandis que sa corrélation avec la vérité augmente.
\item À l'optimum, $\alpha - \gamma\Sigma w^* = \sum_k g_k$, et $\alpha^\top\Sigma^{-1}\alpha - \gamma^2 w^{*\top}\Sigma w^* = \sum_k g_k^\top\Sigma^{-1}(2\alpha - G)$.
\item \textbf{Résultat nommé.} La limite de rotation absorbe 44.0~\% du $\mathrm{IR}^2$ ex ante non contraint, la liquidité 39.4~\%, les neutralités 7.0~\%, les \omterm{def:rs:portfolio-construction-i:constraints}{limites par titre}
3.1~\%, la limite brute et la \omterm{def:rs:portfolio-construction-i:constraints}{neutralité en montant} ne coûtent rien. Après frais, la limite de rotation rapporte le plus~: le ratio net passe de 0.46 à 1.77.
\item Aucune des deux n'est saturée~: l'exposition brute du portefeuille vaut environ 1 une fois que la liquidité et la rotation interviennent, et la \omterm{def:rs:portfolio-construction-i:constraints}{neutralité en montant} ne coûte rien avec un alpha propre à chaque action.
\item Le portefeuille entièrement contraint, présenté avec sa prévision et sa volatilité réalisée, ses frais, les prix fictifs et la décomposition, ainsi que l'erreur type de
son ratio d'information.
\item Quand le nombre de périodes est plusieurs fois celui des titres, ou comme entrée d'un rétrécissement, jamais tel quel avec plus de titres que d'observations.
\item Qu'ils réduisent les poids là où l'erreur d'estimation est la plus grande, si bien qu'une contrainte fausse sur le papier peut réduire le risque réalisé.
\item Sous forme factorielle ($y = X^\top w$), avec un solveur creux ou à séparation d'opérateurs, des démarrages à chaud, et un filtrage des titres qui ne peuvent pas payer leurs coûts.
\item Autant que ses entrées sont justes~: rien si celles-ci sont une covariance empirique de plus de titres que de périodes.
\end{enumerate}
\end{solution}

\begin{solution}{iq:rs:portfolio-construction-i:1}
Ce sont des maximisateurs d'erreurs (Michaud)~: l'optimiseur privilégie les actifs dont les rendements sont surestimés, les variances sous-estimées et les corrélations sous-estimées, qui sont
des erreurs d'estimation. Remèdes~: mettre les alphas à l'échelle du risque, recourir à des modèles factoriels de risque et à des rétrécissements, contraindre, et pénaliser la rotation.
\end{solution}

\begin{solution}{iq:rs:portfolio-construction-i:2}
Le taux auquel l'objectif s'améliorerait par unité d'assouplissement d'une contrainte. On l'utilise pour voir quelles contraintes sont saturées et ce qu'elles coûtent, pour négocier des limites
avec les responsables des risques et, sous forme de vecteurs, pour décomposer ce que chaque contrainte prélève sur le ratio d'information de la prévision.
\end{solution}

\begin{solution}{iq:rs:portfolio-construction-i:3}
La part de l'alpha située le long des expositions neutralisées est perdue~; si l'alpha est propre à chaque action (non corrélé aux expositions), le coût est proche de zéro
(ici de 1.64 à 1.59 ex ante). Jacobs, Levy et Starer ont montré qu'une telle neutralité est en général sous-optimale~; on la retient pour des raisons de risque et de mandat.
\end{solution}

\begin{solution}{iq:rs:portfolio-construction-i:4}
Exprimée en fraction de la valeur quotidienne moyenne échangée rapportée au capital, de sorte que le portefeuille puisse être liquidé en quelques jours~; à mesure que le capital croît, les bornes se resserrent et davantage de titres sont saturés
(28 sur 100 ici à 1 milliard de dollars), le portefeuille se tourne vers les titres liquides et le ratio ex ante baisse~: capacité (chapitre~28).
\end{solution}

\begin{solution}{iq:rs:portfolio-construction-i:5}
La corrélation entre les poids ajustés du risque du portefeuille et la prévision ajustée du risque (Clarke, de Silva et Thorley)~; elle mesure la part de la prévision qui parvient au
portefeuille. Quand la prévision est bruitée, un portefeuille qui ne la suit pas délibérément (une limite de rotation) a un coefficient plus faible et un meilleur ratio réalisé.
\end{solution}

\begin{solution}{iq:rs:portfolio-construction-i:6}
Forme factorielle, avec les expositions comme variables~; un solveur creux de points intérieurs ou de séparation d'opérateurs~; démarrages à chaud à partir de la solution de la veille~; covariance factorielle,
variances spécifiques, expositions et bornes de liquidité précalculées~; filtrage des titres dont l'alpha ne peut couvrir leurs coûts~; contrôles des contraintes saturées et des
prix fictifs avant toute transaction.
\end{solution}
